\begin{abstract}
Scaling each EEG epoch by its own mean and standard deviation is a common preprocessing step in decoding pipelines. We prove that for a baseline-corrected epoch this step maps the pre-stimulus mean onto a sign-flipped, rescaled copy of the post-stimulus mean, so any causal encoder receives post-stimulus information before the stimulus occurs. We derive an averaging law under which the discriminability of the leaked feature grows as the square root of the number of averaged trials, which predicts a small effect in single-trial decoding and a large effect in decoding of averaged, participant-level data. On 211 participants of an affective-picture ERP cohort, pre-stimulus channel means decode three affective conditions at 0.542 balanced accuracy under per-epoch z-scoring against 0.346 under pre-stimulus scaling (chance 0.333). A spiking reservoir from the ARSPI-Net framework, simulated only up to stimulus onset, decodes at 0.420, and the published operating point of that reservoir, 0.487, falls to 0.294 under leak-free scaling, a result reproduced by an independent implementation. EEGNet trained on pre-stimulus samples alone reaches \EEGNETPRE. On five paradigms of the public ERP CORE dataset, the across-participant leak reaches 0.573 to 0.715 in four paradigms and is weakest in the mismatch negativity, and the square-root law predicts the within-participant leak from single-trial accuracy alone. Three of ten widely used EEG decoding codebases scale each sample by its own statistics, and 70\% of 40 recent decoding papers do not report their scaling. Scaling with pre-stimulus or training statistics and reporting pre-stimulus decoding as a null check remove the leak.
\end{abstract}

\begin{IEEEkeywords}
EEG decoding, event-related potentials, normalization, data leakage, spiking reservoir, ARSPI-Net, evaluation methodology.
\end{IEEEkeywords}
